\section{A failure-oriented evaluation protocol}\label{sec:C13:evaluation}
A proposed session evaluation begins with a set of invariants rather than a single completion score. Examples include preserving the current goal, retaining unresolved obligations, keeping authority current, avoiding duplicate external effects, and making every confirmed effect traceable to an authorized intent. Each invariant should identify the component responsible for enforcement and the evidence required to check it. An invariant that depends entirely on a model judge should be labeled accordingly.

Create interruption points before and after each event append, checkpoint commit, model response, tool dispatch, and effect receipt. Inject failures that lose local process memory, delay messages, duplicate deliveries, and reorder independent observations. Use deterministic tool fixtures where the objective is to isolate session mechanics. Then repeat selected cases against realistic external services or environment snapshots to test whether fixture assumptions hold. Both stages matter: the first improves diagnosis, while the second tests external validity.

A current managed-runtime boundary makes this distinction testable. OpenAI's Agents API documentation separates live events from saved items: reconnecting restores saved work rather than replaying every missed event. It also warns that an idle session or a completed turn does not prove every tool succeeded. Required input is read from the current session, not reconstructed solely from historical items~\cite{cumb2026openaievents}. A reconnect test should therefore reconcile item identifiers, current pending requests and effect evidence separately; retransmitting a prior action because a stream ended is not a valid recovery rule.

Metrics should separate restoration correctness, resumed-task success, duplicated effects, lost obligations, reconciliation frequency, and recovery latency. A fast recovery that repeats an irreversible action is not operationally successful. A correct recovery that waits for external reconciliation may be slower but necessary. Report the rate of indeterminate outcomes rather than removing them from the denominator. This prevents an implementation from appearing reliable by declining to classify its hardest failures.

The baseline cohort can include restart from the initial task, restore from a conversation-only checkpoint, restore from typed state plus events, and restore with effect reconciliation. Hold the policy and tools fixed where possible. The comparison asks which information and contracts are required for continuity. It should not be confounded by giving one baseline a stronger model or more generous retries. State exactly which external effects are mockable and which cannot be rolled back.

For concurrency, use a schedule matrix that forces stale reads, conflicting updates, and mid-flight authority changes. Verify that optimistic conflicts are detected and that accepted merges preserve the intended invariants. Storage-level correctness is necessary but insufficient: the semantic content of a memory update can be wrong even when its transaction commits correctly. Combine deterministic checks on identifiers and versions with targeted semantic judgments whose disagreement is reported.

These experiments are proposed and have not been run here. Existing persistence interfaces establish that checkpoints and durable execution can be implemented; they do not settle the workload-specific question of what must be restored. The scientific contribution of a session study is strongest when it connects a particular invariant, an injected failure, and an observed recovery mechanism. An average success score without that chain is weak evidence for durable agent behavior.

